% Part: lambda-calculus
% Chapter: introduction
% Section: basic-pr-lambda

\documentclass[../../../include/open-logic-section]{subfiles}

\begin{document}

\olfileid{lam}{int}{bas}
\olsection{মৌলিক আদিম পুনরাবৃত্ত অপেক্ষকগুলি \usetoken{S}{lambda definable}}

\begin{lem}
$\Zero$, $\Succ$, এবং $\Proj{n}{i}$ অপেক্ষকগুলি !!{lambda definable}।
\end{lem}

\begin{proof}
$\Zero$ স্রেফ
$\lambd[x][\lambd[y][y]]$।

উত্তরসূরি অপেক্ষক~$\Succ$-কে
$\fn{Succ}(u) = \lambd[x][\lambd[y][x(uxy)]]$ দ্বারা সংজ্ঞায়িত করা হয়।
কেন এটি কাজ করে, তা ভেবে দেখো: প্রতিটি সংখ্যাপদ $\num{n}$-কে পুনরাবর্তক
হিসেবে এবং যেকোনো অপেক্ষক~$f$-কে নিলে, $\fn{Succ}(\num{n},f)$ এমন একটি
অপেক্ষক যা নিবেশ~$y$ পেলে~$y$ থেকে শুরু করে $f$-কে $n$ বার প্রয়োগ করে,
তারপর আরও একবার প্রয়োগ করে।

অভিক্ষেপগুলি নিয়ে আর কিছু বলার নেই: $\fn{Proj}^n_i(x_0, \dots,
x_{n-1}) = x_i$। অন্যভাবে বললে, আমাদের প্রথা অনুসারে $\fn{Proj}^n_i$
হলো ল্যাম্বডা পদ $\lambd[x_0][\dots \lambd[x_{n-1}][x_i]]$।
\end{proof}

\end{document}
